% !TEX root = Lecture_Koszul.tex
\section{Tensor and quadratic algebras}

% TODO (ChangeProposal9, Lecture 2 follow-up): Introduce tensor products over
% arbitrary algebras, with the defining identities explicit. References now
% use prop:vector-tensor-universal for the universal property and
% def:vector-tensor for the basis construction. Keep routine
% verification supplementary; retain the existing editorial comments below.

\subsection{Tensor products over an algebra}

Recall the tensor product $M\otimes_A N$ from Definition~\ref{def:module-tensor}, for a right $A$-module $M$ and a left $A$-module $N$.
Its defining identity is $ma\otimes_A n=m\otimes_A an$ for all $m\in M$, $n\in N$, and $a\in A$.

%No semisimplicity or finiteness assumption on $S,M,N$ is needed for the following properties.

\begin{lemma}[Basic properties of tensor products]\label{lem:tensor-properties}
Let $A$ be any $\Bbbk$-algebra, $M$ a right $A$-module, and $N$ a left $A$-module.
\begin{enumerate}
\item For a vector space $W$ and a bilinear map $b:M\times N\to W$ satisfying $b(ma,n)=b(m,an)$ for all $m\in M$, $a\in A$, and $n\in N$, there is a unique linear factorisation (commutative diagram)
\[
\xymatrix@C=75pt@R=22pt{
M\times N \ar[r]^{(m,n)\mapsto m\otimes_A n} \ar[dr]_{\forall \text{ bilinear }b \text{ s.t. } b(-\cdot a,-)=b(-,a\cdot-)\phantom{aaaaaaaaaa}}
& M\otimes_A N \ar@{-->}[d]^{\exists!\,\widetilde b}\\
& W.
}
\]
%Conversely, every linear map $M\otimes_A N\to W$ defines such a bilinear map.
%not sure what this line wants to say really. just ignore it

\item If $B,C$ are $\Bbbk$-algebras and ${}_BM_A$, ${}_AN_C$ are bimodules, then $M\otimes_A N$ is a bimodule ${}_B(M\otimes_A N)_C$ with
\[
b(m\otimes_A n)c=bm\otimes_A nc
\qquad(b\in B,\ c\in C,\ m\in M,\ n\in N).
\]
For right $A$-module maps $u:M\to M'$ and left $A$-module maps $v:N\to N'$, the assignment $(m,n)\mapsto u(m)\otimes v(n)$ induces
\[
u\otimes_A v:M\otimes_A N\longrightarrow M'\otimes_A N',
\qquad m\otimes n\longmapsto u(m)\otimes v(n).
\]
These define functors
\[
(-)\otimes_A N:\Mod A\longrightarrow\Mod C,\qquad
M\otimes_A(-):\Mod(A^\op)\longrightarrow\Mod(B^\op).
\]
On bimodule maps, they retain the outer actions:
\[
(-)\otimes_A N: B\Mod A \longrightarrow B\Mod C,\qquad
M\otimes_A(-):A\Mod C\longrightarrow B\Mod C.
\]
\item For $\Bbbk$-algebras $A,B$, a right $A$-module $U$, a bimodule ${}_AV_B$, and a left $B$-module $W$, there is a natural linear isomorphism
\[
(U\otimes_A V)\otimes_B W\xrightarrow{\sim}U\otimes_A(V\otimes_B W),
\qquad (u\otimes v)\otimes w\longmapsto u\otimes(v\otimes w).
\]
There are also natural module isomorphisms
\[
M\otimes_A A\xrightarrow{\sim}M,\quad m\otimes a\longmapsto ma,
\qquad
A\otimes_A N\xrightarrow{\sim}N,\quad a\otimes n\longmapsto an.
\]
Moreover, these maps are linear with respect to any additional bimodule action on $U, W, M, N$.

\item Tensor products distribute over direct sums in either variable. For arbitrary families $(M_i)_{i\in I}$ of right $A$-modules and $(N_j)_{j\in J}$ of left $A$-modules, indexed by sets $I,J$,
\[
\begin{aligned}
\Bigl(\bigoplus_{i\in I}M_i\Bigr)\otimes_A N
&\xrightarrow{\sim}\bigoplus_{i\in I}(M_i\otimes_A N),
& (m_i)_{i\in I}\otimes n&\longmapsto(m_i\otimes n)_{i\in I},\\
M\otimes_A\Bigl(\bigoplus_{j\in J}N_j\Bigr)
&\xrightarrow{\sim}\bigoplus_{j\in J}(M\otimes_A N_j),
& m\otimes(n_j)_{j\in J}&\longmapsto(m\otimes n_j)_{j\in J}.
\end{aligned}
\]

\item For finite sets $I,J$ and right $A$-modules $M_j,N_i$, there is a natural isomorphism
\begin{equation}\label{eq:hom-finite-sums}
\begin{aligned}
\Hom_A\Bigl(\bigoplus_{j\in J}M_j,\bigoplus_{i\in I}N_i\Bigr)
&\xrightarrow{\sim}\bigoplus_{\substack{i\in I\\j\in J}}\Hom_A(M_j,N_i),\\
f&\longmapsto(\pi_i f\iota_j)_{i\in I,\ j\in J},
\end{aligned}
\end{equation}
where $\iota_j:M_j\to\bigoplus_{j'\in J}M_{j'}$ and $\pi_i:\bigoplus_{i'\in I}N_{i'}\to N_i$ are the inclusions and projections.
\end{enumerate}
\end{lemma}

\begin{supplementary}
\begin{proof}
For (1), Proposition~\ref{prop:vector-tensor-universal} gives a unique linear map $M\otimes_\Bbbk N\to W$, $m\otimes n\mapsto b(m,n)$.
The map vanishes on the defining subspace precisely when $b(ma,n)=b(m,an)$.

For (2), the outer actions preserve that subspace: for $b\in B$, $c\in C$, and $a\in A$,
\[
b(ma\otimes n-m\otimes an)c
=(bm)a\otimes nc-bm\otimes a(nc).
\]
The module axioms follow from those on $M$ and $N$.
For module maps $u,v$,
\[
u(ma)\otimes v(n)=u(m)a\otimes v(n)
=u(m)\otimes av(n)=u(m)\otimes v(an),
\]
so the displayed formula induces $u\otimes_A v$.
The functor identities follow from
$(u'\otimes v')(u(m)\otimes v(n))=u'u(m)\otimes v'v(n)$ and
$(\id_M\otimes\id_N)(m\otimes n)=m\otimes n$.
If $u,v$ preserve the outer actions, then $u\otimes_A v$ does as well.

For (3), both parenthesisations impose
\[
(ua)\otimes v\otimes w=u\otimes(av)\otimes w,\qquad
u\otimes(vb)\otimes w=u\otimes v\otimes(bw)
\quad(a\in A,\ b\in B).
\]
Thus the displayed associativity map and the reverse assignment are well-defined and inverse.
The unit inverses are $M\to M\otimes_A A$, $m\mapsto m\otimes1$, and $N\to A\otimes_A N$, $n\mapsto1\otimes n$.

For (4), the bilinear map $((m_i)_{i\in I},n)\mapsto(m_i\otimes n)_{i\in I}$ satisfies the identity in (1), so it induces the displayed linear map.
For each $i\in I$, let $\iota_i:M_i\to\bigoplus_{j\in I}M_j$ be the inclusion.
The inverse sends $m\otimes n$ in the $i$th summand to $\iota_i(m)\otimes n$ and extends over the direct sum.
Both composites fix every element $m\otimes n$ from a summand.
The other variable is handled in the same way.
\end{proof}
\end{supplementary}


\begin{advanced}
\begin{lemma}[Hom, direct sums, and direct products]\label{lem:hom-sums-products}
Let $M,N$ be right $A$-modules, and let $(M_i)_{i\in I}$ and $(N_i)_{i\in I}$ be arbitrary families of right $A$-modules, indexed by a set $I$.
Write $\iota_i:M_i\to\bigoplus_{j\in I}M_j$ for the inclusions and $\pi_i:\prod_{j\in I}N_j\to N_i$ for the projections.
There are natural linear isomorphisms
\[
\begin{aligned}
\Hom_A\Bigl(\bigoplus_{i\in I}M_i,N\Bigr)&\xrightarrow{\sim}\prod_{i\in I}\Hom_A(M_i,N),
& f&\longmapsto(f\circ\iota_i)_{i\in I},\\
\Hom_A\Bigl(M,\prod_{i\in I}N_i\Bigr)&\xrightarrow{\sim}\prod_{i\in I}\Hom_A(M,N_i),
& f&\longmapsto(\pi_i\circ f)_{i\in I}.
\end{aligned}
\]
There is also a natural injection
\[
\bigoplus_{i\in I}\Hom_A(M,N_i)\longrightarrow
\Hom_A\Bigl(M,\bigoplus_{i\in I}N_i\Bigr),
\qquad (f_i)_{i\in I}\longmapsto\bigl(m\mapsto(f_i(m))_{i\in I}\bigr),
\]
which is an isomorphism when $M$ is finitely generated.
\end{lemma}

\begin{proof}
The inverse of the first map sends $(f_i)_{i\in I}$ to $(m_i)_{i\in I}\mapsto\sum_{i\in I}f_i(m_i)$; the sum is finite because $(m_i)_{i\in I}$ has finite support.
The inverse of the second sends $(f_i)_{i\in I}$ to $m\mapsto(f_i(m))_{i\in I}$; a direct product imposes no support condition.
For the third map, coordinate projections recover each $f_i$, proving injectivity.
If $m_1,\ldots,m_r$ generate $M$, the images of these generators under $f:M\to\bigoplus_{i\in I}N_i$ are supported on a common finite subset $J\subset I$.
Then $f(M)\subset\bigoplus_{i\in J}N_i$, so $f$ comes from a finite family of maps.
All three constructions commute with precomposition and postcomposition.
\end{proof}

\begin{proposition}[Tensor products and direct products]\label{prop:tensor-direct-products}
Let $A$ be a $\Bbbk$-algebra, $M,M_i$ right $A$-modules, and $N,N_i$ left $A$-modules, with $i\in I$.
The coordinate projections define linear maps
\begin{equation}\label{eq:module-tensor-products}
\begin{aligned}
\kappa_N:\Bigl(\prod_iM_i\Bigr)\otimes_A N
&\longrightarrow\prod_i(M_i\otimes_A N),
& (m_i)_i\otimes n&\longmapsto(m_i\otimes n)_i,\\
\lambda_M:M\otimes_A\Bigl(\prod_iN_i\Bigr)
&\longrightarrow\prod_i(M\otimes_A N_i),
& m\otimes(n_i)_i&\longmapsto(m\otimes n_i)_i.
\end{aligned}
\end{equation}
\begin{enumerate}
\item If $I$ is finite, both maps are isomorphisms.
\item The map $\kappa_N$ is an isomorphism if $N$ is a direct summand of $A^{\oplus r}$ as a left module for some finite $r$.
The corresponding condition on the right module $M$ makes $\lambda_M$ an isomorphism.
\item If $A$ is finite-dimensional semisimple, both maps are injective.
Moreover, $\kappa_N$ is an isomorphism when $N$ is finite-dimensional, and $\lambda_M$ is an isomorphism when $M$ is finite-dimensional.
\end{enumerate}
\end{proposition}

\begin{proof}
The formulas in \eqref{eq:module-tensor-products} are bilinear. For $s\in A$, the component identities
$(m_is)\otimes n=m_i\otimes sn$ and $(ms)\otimes n_i=m\otimes sn_i$
show that both formulas give well-defined maps from the indicated tensor products.
For finite $I$, products equal sums, and Lemma~\ref{lem:tensor-properties}(4) proves (1).

For $N=A$, the map $\kappa_A$ is the identity on $\prod_iM_i$ under $m\otimes s\mapsto ms$.
For $N=A^{\oplus r}$, both sides are $(\prod_iM_i)^{\oplus r}$, since a finite direct sum commutes with a direct product.
If $N\oplus N'\cong A^{\oplus r}$, the comparison for $N\oplus N'$ is $\kappa_N\oplus\kappa_{N'}$.
Its bijectivity implies the bijectivity of $\kappa_N$.
This proves the assertion about $\kappa_N$ in (2).

Now let $A$ be finite-dimensional semisimple.
For finite-dimensional $N$, Corollary~\ref{cor:semisimple-summands} gives (2).
For arbitrary $N$, take $z=\sum_{\ell=1}^r m^{(\ell)}\otimes n_\ell$ in $\Ker\kappa_N$ and set $N_0=\sum_{\ell=1}^r An_\ell$.
The submodule $N_0$ is finite-dimensional and has a complement in $N$ by Theorem~\ref{thm:artin-wedderburn}.
Thus the inclusions
\[
\Bigl(\prod_iM_i\Bigr)\otimes_A N_0
\longrightarrow\Bigl(\prod_iM_i\Bigr)\otimes_A N,
\qquad
M_i\otimes_A N_0\longrightarrow M_i\otimes_A N
\]
are injective, by the direct sum formula in Lemma~\ref{lem:tensor-properties}(4).
Regard $z$ as an element with second factors in $N_0$.
The equality $\kappa_N(z)=0$ then gives $\kappa_{N_0}(z)=0$.
Aince $\kappa_{N_0}$ is an isomorphism, $z=0$.

For $\lambda_M$, apply the same arguments over $A^\op$, using
$N\otimes_{A^\op}M\cong M\otimes_A N$, $n\otimes m\mapsto m\otimes n$.
\end{proof}

The finite-dimensional hypothesis on the fixed factor cannot be dropped from the isomorphism assertion: Example~\ref{ex:tensor-product-product-failure} gives a non-surjective comparison even for vector spaces.
Thus tensor products always distribute over direct sums, but not over arbitrary direct products.

\begin{remark}\label{rem:infinite-products}
For infinite families, $\Hom$ need not preserve direct sums in its second variable.
For $M=\bigoplus_{i\geq1}\Bbbk v_i$, the identity $M\to M$ does not come from $\bigoplus_i\Hom_\Bbbk(M,\Bbbk v_i)$: its projection onto every $\Bbbk v_i$ is non-zero.
Thus the injection in Lemma~\ref{lem:hom-sums-products} need not be surjective.

A product in the source of $\Hom$ behaves differently from a product in the target.
For example, a linear map $\prod_{i\geq1}\Bbbk\to\Bbbk$ need not be determined by its restrictions to the coordinate copies of $\Bbbk$.
Indeed, $(1,1,\ldots)\notin\bigoplus_{i\geq1}\Bbbk$.
Define a functional on $\bigoplus_{i\geq1}\Bbbk+\Bbbk(1,1,\ldots)$ to vanish on the direct sum and take value $1$ at $(1,1,\ldots)$, and extend it to the product by extending a basis.
The resulting non-zero functional vanishes on every coordinate copy.
\end{remark}
\end{advanced}



\phantomsection\label{not:tensor-powers}
For a bimodule ${}_AV_A$, define
\[
V^{\otimes_A 0}:=A,\qquad
V^{\otimes_A i}:=
\underbrace{V\otimes_A\cdots\otimes_A V}_{i\text{ factors}}
\quad(i\geq1).
\]
Note that there is no parenthesis in the definition of $V^{\otimes_A i}$ since Lemma~\ref{lem:tensor-properties}(3) says that tenso product is an associatitive operation.

\begin{advanced}
Let $A,B,C$ be $\Bbbk$-algebras.
For bimodules ${}_AM_B$, ${}_CX_B$, we have $\Hom_B(M,X)\in C\Mod A$, with
\begin{equation}\label{eq:general-tensor-dual-actions}
(cfa)(m)=cf(am)
\quad(f\in\Hom_B(M,X),\ c\in C,\ a\in A,\ m\in M).
\end{equation}
If $X\in\Mod B$ has no specified left action, the same formula defines the right $A$-action $(fa)(m)=f(am)$.
The left-module version follows by passing to opposite algebras.

\begin{proposition}[Tensor and Hom]\label{prop:tensor-duals}
% TODO (ChangeProposal9, Lecture 2 follow-up): Explain the general case for
% an (A,B)-bimodule M and a (B,C)-bimodule N. For right duals, require M_B
% finitely generated projective; for left duals, require _B N finitely generated
% projective. Finite k-dimension alone does not suffice over a general algebra.
% Keep the idempotent-matching calculation in the main exposition and move
% routine verification to supplementary material. Use the one-sided bimodule
% actions from eq:bimodule-dual-actions, independently of a symmetrising form.
Consider bimodules ${}_AM_B$, ${}_BN_C$, a right $C$-module $X$, and a left $A$-module $Y$.
There are maps, natural in $M,N$,
\begin{equation}\label{eq:general-tensor-duals}
\begin{aligned}
\theta_r:\Hom_C(N,X)\otimes_B\Hom_B(M,B)
&\longrightarrow\Hom_C(M\otimes_B N,X),\\
\theta_r(g\otimes f)(m\otimes n)&=g(f(m)n),\\
\theta_\ell:\Hom_{B^\op}(N,B)\otimes_B\Hom_{A^\op}(M,Y)
&\longrightarrow\Hom_{A^\op}(M\otimes_B N,Y),\\
\theta_\ell(g\otimes f)(m\otimes n)&=f(mg(n)).
\end{aligned}
\end{equation}
The maps $\theta_r$ and $\theta_\ell$ are right $A$-linear and left $C$-linear, respectively.
\begin{enumerate}
\item The map $\theta_r$ is an isomorphism if $M\in\add B$ as a right $B$-module.
\item The map $\theta_\ell$ is an isomorphism if ${}_BN\in\add_{\mod B^\op}(B)$.
\end{enumerate}
Here $\add B=\add_{\mod B}(B)$, as in Example~\ref{ex:add-category}.
If ${}_CX_C$ and ${}_AY_A$ are bimodules, the respective maps also preserve the additional actions and are morphisms in $C\Mod A$.
\end{proposition}

\begin{proof}
\emph{Well-definedness of $\theta_r$.}
For $f\in\Hom_B(M,B)$, $g\in\Hom_C(N,X)$, and $b\in B$,
\[
g(f(mb)n)=g(f(m)bn),\qquad
(gb)(f(m)n)=g(bf(m)n)=g((bf)(m)n).
\]
These identities respect the defining identifications in both tensor products.
Right $C$-linearity follows from $g(f(m)nc)=g(f(m)n)c$ for $c\in C$.
For $a\in A$,
\[
\theta_r(g\otimes fa)(m\otimes n)
=g(f(am)n)=(\theta_r(g\otimes f)a)(m\otimes n).
\]
If $X$ also has a commuting left $C$-action, replacing $g$ by $cg$ multiplies the value by $c$, so $\theta_r$ preserves that action as well.

\emph{Bijectivity of $\theta_r$.}
Regard $M$ only as a right $B$-module.
For $M=B$, evaluation at $1$ identifies $\Hom_B(B,B)\cong B$, while multiplication identifies $B\otimes_B N\cong N$.
Under these isomorphisms, $\theta_r$ becomes
\[
\Hom_C(N,X)\otimes_B B\longrightarrow\Hom_C(N,X),
\qquad g\otimes b\longmapsto gb,
\]
which is an isomorphism by Lemma~\ref{lem:tensor-properties}(3).
Finite direct sums commute with both constructions by Lemma~\ref{lem:tensor-properties}(4) and \eqref{eq:hom-finite-sums}.
Thus $\theta_r$ is bijective for $M=B^{\oplus r}$, where $r\geq0$ is finite.
If $M\in\add B$, choose $M'$ with $M\oplus M'\cong B^{\oplus r}$.
The comparison for $M\oplus M'$ is the direct sum of the comparisons for $M$ and $M'$, so both summand maps are bijective.

\emph{The left comparison.}
For $f\in\Hom_{A^\op}(M,Y)$ and $g\in\Hom_{B^\op}(N,B)$, well-definedness follows from
\[
f(mb\,g(n))=f(mg(bn)),\qquad
f(m(gb)(n))=f(mg(n)b)=(bf)(mg(n))
\quad(b\in B).
\]
Left $A$-linearity of $f$ makes the resulting functional left $A$-linear.
For $c\in C$,
\[
\theta_\ell(cg\otimes f)(m\otimes n)
=f(mg(nc))=(c\theta_\ell(g\otimes f))(m\otimes n).
\]
If $Y$ also has a commuting right $A$-action, replacing $f$ by $fa$ multiplies the value by $a$, so $\theta_\ell$ preserves that action as well.
For $N=B$ as a left module, evaluation at $1$ identifies $\Hom_{B^\op}(B,B)\cong B$ and turns $\theta_\ell$ into
\[
B\otimes_B\Hom_{A^\op}(M,Y)\longrightarrow\Hom_{A^\op}(M,Y),
\qquad b\otimes f\longmapsto bf.
\]
The same finite-sum and direct-summand argument proves (2).

\emph{Naturality.}
Let $p:M\to M'$ and $q:N\to N'$ be bimodule maps.
For $f\in\Hom_B(M',B)$ and $g\in\Hom_C(N',X)$, applying $\theta_r$ before or after precomposition gives $g(f(p(m))q(n))$ at $m\otimes n$.
For $f\in\Hom_{A^\op}(M',Y)$ and $g\in\Hom_{B^\op}(N',B)$, the two orders give $f(p(m)g(q(n)))$ for $\theta_\ell$.
Thus both comparisons are natural.
\end{proof}
\end{advanced}

Let $S$ be finite-dimensional semisimple, and consider bimodules ${}_SV_S$, ${}_SW_S$ finitely generated on the right and left, respectively.
Corollary~\ref{cor:semisimple-summands} gives $V_S\in\add S$ and ${}_SW\in\add_{\mod S^\op}(S)$.
Thus the semisimple case of tensor duality, Proposition~\ref{prop:tensor-duals}, gives natural bimodule isomorphisms \comm{you always call this tensor-dual isomorphism, but this name is very ambiguous... it is more like anti-distributivity of $S$-dual over tensor-over-$S$.  you can use emphasise on this phrase to make it more easy to remember, if you need to use word, then say ``we refer to this as anti-distributivity of $S$-dual" from now on; then apply changes for any further use of this (whenever you say `tensor-dual isom', it should be clarified)}
\begin{equation}\label{eq:tensor-duals}
\theta_r:W^*\otimes_S V^*\xrightarrow{\sim}(V\otimes_S W)^*,
\qquad
\theta_\ell:{}^*W\otimes_S{}^*V\xrightarrow{\sim}{}^*(V\otimes_S W),
\end{equation}
\begin{equation}\label{eq:tensor-evaluations}
\theta_r(g\otimes f)(v\otimes w)=g(f(v)w),
\qquad
\theta_\ell(g\otimes f)(v\otimes w)=f(vg(w)).
\end{equation}
For $M\in\{V,W\}$, the actions from \eqref{eq:bimodule-dual-actions} are
\begin{equation}\label{eq:tensor-dual-actions}
(sft)(m)=sf(tm),\qquad (sgt)(m)=g(ms)t
\quad\text{for all }s,t\in S,\ f\in M^*,\ g\in{}^*M,\ m\in M.
\end{equation}

\begin{hidden}
\begin{convention}[Tensor duals]
For $S,V,W$ as in \eqref{eq:tensor-duals}, identify the dual of a tensor product with the tensor product of the corresponding duals in reverse order, using the maps in \eqref{eq:tensor-evaluations}.
No basis or symmetrising form is needed.
Comparisons of $M^*$, $ {}^*M$, and $D(M)$ instead use a chosen symmetrising form, as in \eqref{eq:dual-comparison}.
\end{convention}
\end{hidden}

\begin{advanced}
\begin{example}\label{ex:tensor-dual-finite-generation}
Finite generation alone does not suffice in Proposition~\ref{prop:tensor-duals} over a general algebra.
Take $A=C=\Bbbk$, $B=\Bbbk[x]$, and $X=Y=\Bbbk$.
Consider bimodules ${}_\Bbbk M_B$, ${}_BN_\Bbbk$ with $M=N=B/(x)\cong\Bbbk$ and the outer $\Bbbk$-actions given by scalar multiplication.
Every $f\in\Hom_B(M,B)$ satisfies $f(\overline1)x=f(\overline1x)=0$, hence $f=0$.
Consequently
\[
\Hom_\Bbbk(N,\Bbbk)\otimes_B\Hom_B(M,B)=0,\qquad
\Hom_\Bbbk(M\otimes_B N,\Bbbk)\cong\Bbbk.
\]
Both $M_B$ and $ {}_B N$ are generated by $\overline1$, so neither direct-summand hypothesis can be replaced by finite generation alone.
For the left comparison, $\Hom_{B^\op}(N,B)=0$, whereas $\Hom_\Bbbk(M\otimes_B N,\Bbbk)\cong\Bbbk$.
\end{example}

For $S,V,W$ as in \eqref{eq:tensor-duals}, take $A=B=S$, $C=\Bbbk$, $M=V$, $N=W$, and $X=\Bbbk$ in Proposition~\ref{prop:tensor-duals}(1).
Since $\Hom_\Bbbk(W,\Bbbk)=D(W)$, we obtain the natural bimodule isomorphism \cite[\S2.1]{GI20}
\[
\begin{aligned}
\theta_D:D(W)\otimes_S V^*&\xrightarrow{\sim}D(V\otimes_S W),\\
\theta_D(\lambda\otimes f)(v\otimes w)&=\lambda(f(v)w).
\end{aligned}
\]
The actions on $D(W)$ are those in \eqref{eq:bimodule-dual-actions}.
For $s\in S$, naturality with respect to $W\to W$, $w\mapsto ws$, gives left $S$-linearity of $\theta_D$.

Suppose now that $S$ is equipped with a symmetrising form $\tau:S\to\Bbbk$, as in Definition~\ref{def:symmetrising-form}.
The two-sided dual of Definition~\ref{def:en-dual} satisfies
\begin{equation}\label{eq:tensor-two-sided-dual}
\begin{aligned}
\theta_{\en}:W^*\otimes_S V^{*_{\en}}
&\xrightarrow{\sim}(V\otimes_S W)^{*_{\en}},\\
\theta_{\en}(g\otimes F)(v\otimes w)
&=\sum_j g(a_jw)\otimes b_j^\op,
\qquad F(v)=\sum_j a_j\otimes b_j^\op.
\end{aligned}
\end{equation}
We check that the formula defines a map on the two indicated tensor products and preserves both outer actions.
For example, $F(vs)=F(v)(s\otimes1)$ replaces $a_j$ by $a_js$, while
$(Fs)(v)=(1\otimes s^\op)F(v)$ replaces $b_j$ by $b_js$.
These give, respectively, compatibility with $vs\otimes w=v\otimes sw$ and with the right outer action.
Compatibility with $(gs)\otimes F=g\otimes(sF)$ follows from $(gs)(a_jw)=g(sa_jw)$.
The other actions follow from right $S$-linearity of $g$ and $F(tv)=F(v)(1\otimes t^\op)$.

The comparison maps $\gamma_M$ in \eqref{eq:en-dual-comparison} give
\[
\gamma_{V\otimes_S W}\circ\theta_{\en}
=\theta_r\circ(\id_{W^*}\otimes\gamma_V),
\]
since both sides send $g\otimes F$, evaluated at $v\otimes w$, to $\sum_j g(a_jw)\tau(b_j)$.
Thus $\theta_{\en}$ is an isomorphism by Theorem~\ref{prop:dual-comparison}.
Precomposition shows that $\theta_D$ and $\theta_{\en}$ are natural.
No assumption that $S^\en$ is semisimple is used; compare Remark~\ref{rem:enveloping-semisimple}.

For the maps $\alpha_M(f)=\tau\circ f$ in \eqref{eq:dual-comparison},
\[
\alpha_{V\otimes_S W}\circ\theta_r
=\theta_D\circ(\alpha_W\otimes\id_{V^*}),
\]
because both sides send $g\otimes f$ to the functional $v\otimes w\mapsto\tau(g(f(v)w))$.
\end{advanced}


\subsection{Generated algebras}

\begin{definition}
Let $A$ be graded with $S=A_0$.
Multiplication makes $A_1$ an $S$-bimodule and gives
\begin{equation}\label{eq:multiplication-maps}
\mu_i:A_1^{\otimes_S i}\longrightarrow A_i,\qquad
a_1\otimes\cdots\otimes a_i\longmapsto a_1\cdots a_i
\quad(i\geq 1),\qquad \mu_0=\id_S.
\end{equation}
The algebra $A$ is \defn{generated in degree $1$ (over $S$)} if every $\mu_i$ in \eqref{eq:multiplication-maps} is surjective.
\end{definition}

Equivalently, $S$ and $A_1$ generate $A$ as an algebra.

\begin{example}
The algebra $\Bbbk A_n$ of Example~\ref{ex:triangular-grading} is generated in degree $1$ over $S=\bigoplus_i\Bbbk e_{ii}$.
%\[
%e_{ij}=e_{i,i+1}e_{i+1,i+2}\cdots e_{j-1,j}\quad(1\leq i<j\leq n).
%\]
%For the grading $\deg e_{ij}=2(j-i)$, the same algebra has $A_1=0$ and $A_2\neq 0$, so it is not generated in degree $1$.
\end{example}

\begin{example}\label{ex:polynomial-generators}
The polynomial algebra $\Bbbk[x_1,\ldots,x_n]$, with each $\deg x_i=1$ for all $i$, is generated in degree $1$ over $\Bbbk$.
Indeed, every monomial of total degree $d$ is a product of $d$ variables, and these monomials span the degree $d$ component.
\end{example}

\begin{supplementary}
\begin{example}\label{ex:triangular-polynomials}
\emergencystretch=1em
Let
\[
A=\begin{pmatrix}\Bbbk[x]&\Bbbk[x]\\0&\Bbbk[x]\end{pmatrix},
\qquad \deg(x^m e_{ij})=m+j-i
\quad(m\geq0,\ 1\leq i\leq j\leq2).
\]
Then
\[
A_0=\Bbbk e_{11}\oplus\Bbbk e_{22},\qquad
A_d=\Bbbk x^d e_{11}\oplus\Bbbk x^d e_{22}\oplus\Bbbk x^{d-1}e_{12}
\quad(d\geq1).
\]
It is generated in degree $1$ over $A_0$, since
\[
x^d e_{11}=(xe_{11})^d,\qquad
x^d e_{22}=(xe_{22})^d,\qquad
x^{d-1}e_{12}=(xe_{11})^{d-1}e_{12}
\quad(d\geq1).
\]
Here $A_0$ has two simple summands, $A_d\neq0$ for every $d\geq0$, and $A$ is non-commutative: $e_{11}e_{12}=e_{12}$ whereas $e_{12}e_{11}=0$.
\end{example}
\end{supplementary}

\subsection{Tensor algebras and the universal property}

% TODO (ChangeProposal9, Lecture 2 follow-up): Keep the tensor-algebra universal
% property and its explicit formula in the main text; shorten its proof or put
% routine verification in supplementary material. Address the remaining local
% editorial comments when revising Lecture 2, not as part of the Lecture 1 edit.

\begin{definition}\label{def:tensor-algebra}
Let $S$ be a $\Bbbk$-algebra and consider a bimodule ${}_SV_S$.
The \defn{tensor algebra of $V$ over $S$} is
\[
T_S(V):=\bigoplus_{i\in\Z}T_S(V)_i,
\quad\text{where}\quad
T_S(V)_i=
\begin{cases}
V^{\otimes_S i},&i\geq1,\\
S,&i=0,\\
0,&i<0,
\end{cases}.
\]
For $v_1,\ldots,v_{p+q}\in V$, multiplication is determined $\Bbbk$-bilinearly by
\[
(v_1\otimes_S\cdots\otimes_S v_p)(v_{p+1}\otimes_S\cdots\otimes_S v_{p+q})
=v_1\otimes_S\cdots\otimes_S v_{p+q}
\qquad(p,q\geq1),
\]
and by the left and right $S$-actions when a factor has degree $0$.
\end{definition}

We also call $T_S(V)$ the tensor algebra generated by $V$ over $S$.
Every element has finitely many non-zero homogeneous components.
The unit is $1_S$, and associativity follows from Lemma~\ref{lem:tensor-properties}(3).
Write $v_1\cdots v_d$ for $v_1\otimes_S\cdots\otimes_S v_d$ in $T_S(V)$.
In a quotient algebra (including $T_S(V)$ itself) of $T_S(V)$, the notation $v_1\cdots v_d$ denotes the image of $v_1\otimes_S\cdots\otimes_S v_d$.

\begin{example}[Polynomial and exterior algebras]\label{ex:polynomial-exterior}
Let $V,W$ be finite-dimensional vector spaces with bases $\{x_i\mid1\leq i\leq n\}$ and $\{y_j\mid1\leq j\leq m\}$, respectively.
The \defn{symmetric algebra} of $V$ and the \defn{exterior algebra} of $W$ are
\begin{equation}\label{eq:polynomial-exterior}
\begin{aligned}
\SymAlg(V)&=T_\Bbbk(V)/(x_ix_j-x_jx_i\mid1\leq i<j\leq n)
            \cong\Bbbk[x_1,\ldots,x_n],\\
\bigwedge W&=T_\Bbbk(W)/(w^2\mid w\in W) = T_\Bbbk(W)/(y_j^2, y_iy_j+y_jy_i \mid 1\le i<j\le m).
\end{aligned}
\end{equation}
They are generated in degree $1$, with
\[
\SymAlg(V)_0=\Bbbk=(\bigwedge W)_0,\qquad
\SymAlg(V)_1=V,\qquad(\bigwedge W)_1=W.
\]
Their respective bases are
\[
\begin{gathered}
\{x_1^{a_1}\cdots x_n^{a_n}\mid a_1,\ldots,a_n\in\Z,\ a_1,\ldots,a_n\geq0\},\\
\{y_{j_1}\cdots y_{j_d}\mid0\leq d\leq m,\ 1\leq j_1<\cdots<j_d\leq m\}.
\end{gathered}
\]
\end{example}

\begin{supplementary}
The bases in Example~\ref{ex:polynomial-exterior} can be checked directly.
For $\SymAlg(V)$, commuting adjacent generators gives the ordered monomials, which are independent under the map to the polynomial ring.
For $\bigwedge W$, the relations reduce every product of basis elements to a multiple of a strictly increasing product or to zero.
For $1\leq d\leq m$, fix $j_1<\cdots<j_d$ and consider the linear functional $W^{\otimes_\Bbbk d}\to\Bbbk$ sending $w_1\otimes\cdots\otimes w_d$ to the determinant of their coordinates in rows $j_1,\ldots,j_d$ relative to the basis $\{y_j\mid1\leq j\leq m\}$.
The determinant vanishes whenever two adjacent factors agree, hence on the degree $d$ component of $(w^2\mid w\in W)$.
On increasing products it is $1$ on $y_{j_1}\cdots y_{j_d}$ and $0$ on the others, proving independence in every characteristic.
\end{supplementary}

\begin{lemma}\label{lem:idempotent-tensors}
Let $S=\bigoplus_{i=1}^r\Bbbk e_i$ with coordinate idempotents, and consider bimodules ${}_SV_S$, ${}_SW_S$.
There is a bimodule isomorphism
\begin{equation}\label{eq:idempotent-tensor-decomposition}
V\otimes_S W\cong\bigoplus_{i,j,k=1}^r e_iVe_j\otimes_\Bbbk e_jWe_k.
\end{equation}
Consequently, after choosing a basis of each $e_iVe_j$, a basis of $V^{\otimes_S d}$ for $d\geq1$ consists of products of $d$ chosen basis elements whose adjacent idempotents match.
\end{lemma}
\begin{supplementary}
\begin{proof}
The orthogonal idempotents sum to $1$, so
\begin{equation}\label{eq:composable-tensors}
V=\bigoplus_{i,j=1}^r e_iVe_j,\qquad
(e_iVe_j)\otimes_S(e_pWe_k)=0
\quad(1\leq i,j,p,k\leq r,\ j\neq p).
\end{equation}
Indeed, $v\otimes_S w=ve_j\otimes_S w=v\otimes_S e_jw=0$ for $v\in e_iVe_j$, $w\in e_pWe_k$, and $j\neq p$.
The map from right to left in \eqref{eq:idempotent-tensor-decomposition} sends $v\otimes_\Bbbk w$ to $v\otimes_S w$ on each summand.
Its inverse is
\[
V\otimes_S W\longrightarrow\bigoplus_{i,j,k=1}^r e_iVe_j\otimes_\Bbbk e_jWe_k,
\qquad
v\otimes_S w\longmapsto\sum_{i,j,k=1}^r e_ive_j\otimes_\Bbbk e_jwe_k.
\]
Checking that these two assignments induce well-defined, mutually inverse maps is left to the reader.
Iterating \eqref{eq:idempotent-tensor-decomposition} and using Definition~\ref{def:vector-tensor} gives the stated basis.
\end{proof}
\end{supplementary}


\begin{example}\label{ex:tensor-A3}
We apply Lemma~\ref{lem:idempotent-tensors} to realise $\Bbbk A_3$ as a tensor algebra.
Let $S=\Bbbk^{\oplus3}$ with coordinate idempotents $e_1,e_2,e_3$, and take the bimodule ${}_SV_S$ with $V=\Bbbk a\oplus\Bbbk b$, $e_1ae_2=a$, and $e_2be_3=b$; all other idempotent actions are zero.
By Lemma~\ref{lem:idempotent-tensors},
\[
V\otimes_S V=\Bbbk(a\otimes b),\qquad
V^{\otimes_S3}=0,\qquad
T_S(V)=S\oplus\Bbbk a\oplus\Bbbk b\oplus\Bbbk ab.
\]
Here $ab=a\otimes b$.
The linear map $T_S(V)\to\Bbbk A_3$ given by $e_i\mapsto e_{ii}$, $a\mapsto e_{12}$, $b\mapsto e_{23}$, and $ab\mapsto e_{13}$ preserves multiplication and sends a basis to a basis, hence is an algebra isomorphism.
\end{example}

\comm{at some point before universal property prop, it should be clarified that tensor algebra is a graded algebra.  the current definition just says something of degree 0, without saying degree 0 in where.  If necessary, we will refer to the canonical grading structure $\deg S=0, \deg V=1$ as tensor grading.}

\begin{proposition}[Universal property of the tensor algebra]\label{prop:tensor-universal}
Let $S$ be a $\Bbbk$-algebra, consider a bimodule ${}_SV_S$, and let $\iota_S:S\to T_S(V)$ and $\iota_V:V\to T_S(V)$ be the canonical inclusions.
Given a unital $\Bbbk$-algebra homomorphism $\varphi:S\to B$ and a bimodule map $f:{}_SV_S\to{}_SB_S$, where $sbt=\varphi(s)b\varphi(t)$ on $B$, there is a unique $\Bbbk$-algebra homomorphism $\widetilde f:T_S(V)\to B$ making the diagram commute:
\begin{equation}\label{eq:tensor-universal}
\xymatrix@C=42pt@R=28pt{
V \ar@{^{(}->}[r]^{\iota_V} \ar[dr]_{\forall\,\text{bimodule map }f\phantom{...}}
& T_S(V) \ar@{-->}[d]^{\exists!\,\widetilde f}
& S \ar@{_{(}->}[l]_{\iota_S} \ar[dl]^{\phantom{...}\forall\,\text{algebra map }\varphi}\\
& B. &
}
\end{equation}
\end{proposition}
\begin{supplementary}
\begin{proof}
Define $\widetilde f$ on $S$ and each tensor power by
\begin{equation}\label{eq:tensor-extension}
\widetilde f(s)=\varphi(s),\qquad
\widetilde f(v_1\otimes_S\cdots\otimes_S v_i)=f(v_1)\cdots f(v_i)
\quad(s\in S,\ i\geq1).
\end{equation}
The identity $f(vs)f(w)=f(v)\varphi(s)f(w)=f(v)f(sw)$ shows that replacing adjacent factors $vs\otimes w$ by $v\otimes sw$ does not change the value in \eqref{eq:tensor-extension}.
The formulas therefore induce linear maps on the tensor powers, which extend to their direct sum.
Multiplication is preserved by \eqref{eq:tensor-extension}, including the degree $0$ factors because $f(sv)=\varphi(s)f(v)$ and $f(vs)=f(v)\varphi(s)$.
Also, $\widetilde f(1_S)=1_B$.
Every algebra homomorphism making \eqref{eq:tensor-universal} commute must have the values in \eqref{eq:tensor-extension}, proving uniqueness.
\end{proof}

If $B$ is graded, $\varphi(S)\subset B_0$, and $f(V)\subset B_1$, then \eqref{eq:tensor-extension} gives $\widetilde f(V^{\otimes_S i})\subset B_i$ for every $i\geq0$.
Thus $\widetilde f$ is a graded algebra homomorphism.
\end{supplementary}


\begin{corollary}\label{cor:tensor-presentation}
If $A$ is generated in degree $1$ over $S=A_0$, then the multiplication map defines a surjective graded algebra homomorphism
\[
\mu:T_S(A_1)\onto A,\qquad \mu|_S=\id_S, \qquad 
a_1\otimes_S\cdots\otimes_S a_i\longmapsto a_1\cdots a_i
\quad \forall i\geq1.
\]
Its degree $i$ component is $\mu_i$ from \eqref{eq:multiplication-maps}.
\end{corollary}

\begin{supplementary}
\begin{proof}
In Proposition~\ref{prop:tensor-universal}, take $V=A_1$, $B=A$, and the inclusions
\[
\varphi:S\hookrightarrow A,\quad s\longmapsto s,
\qquad f:A_1\hookrightarrow A,\quad a\longmapsto a.
\]
Then \eqref{eq:tensor-extension} gives $\widetilde f=\mu$.
The map $\mu$ preserves degrees, and each component $\mu_i$ is surjective by the assumption that $A$ is generated in degree $1$.
\end{proof}
\end{supplementary}

\begin{definition}\label{def:relations}
For the map $\mu$ in Corollary~\ref{cor:tensor-presentation}, the homogeneous two-sided ideal $I:=\Ker\mu$ is the \defn{ideal of relations}.
\end{definition}

\comm{this paragraph, up to `Homegentiy is essential...', should be inside the corollary??}
Note that $I_0=I_1=0$ since $\mu_0=\id_S$ and $\mu_1=\id_{V=A_1}$.
Since $\mu$ preserves degrees, its kernel is homogeneous, and
\[
T_S(V)/I\xrightarrow{\sim}A,\qquad z+I\longmapsto\mu(z)
\]
is a graded algebra isomorphism.
The bimodule $V$ specifies the generators and their left and right $S$-actions; $I$ specifies the relations between their products.
Homogeneity is essential for the quotient grading: the relation $x-x^2$ in Example~\ref{ex:nonhomogeneous-ideal} does not satisfy this condition.

\subsection{Quadratic algebras}

We now consider algebras whose ideal of relations is generated in degree $2$.

\begin{definition}\label{def:quadratic} 
A graded algebra $A$ is \defn{quadratic} if there exist a finite-dimensional semisimple $\Bbbk$-algebra $S$, a finite-dimensional $S$-bimodule $V$, and an $S$-sub-bimodule $R\subset V\otimes_S V$ such that $A\cong T_S(V)/(R)$ as graded algebras, where $(R)$ is the two-sided ideal of $T_S(V)$ generated by $R$.
\end{definition}

An isomorphism $A\cong T_S(V)/(R)$ as in Definition~\ref{def:quadratic} is called a \defn{quadratic presentation} of $A$.

\begin{proposition}\label{prop:quadratic-basic}
A quadratic algebra $A=T_S(V)/(R)$ has $A_0=S$, $A_1=V$, and $A_2=(V\otimes_S V)/R$, is locally finite-dimensional, and is generated in degree $1$ over $S$.
\end{proposition}

\begin{proof}
The ideal $(R)$ has no components in degrees $0$ and $1$.
Its degree $2$ component is exactly $R$.
Each $A_i$ is a quotient of the finite-dimensional space $V^{\otimes_S i}$, and multiplication maps $V^{\otimes_S i}$ onto $A_i$.
\end{proof}

Thus a quadratic presentation recovers $S\cong A_0$, $V\cong A_1$, and $R\cong\Ker(A_1\otimes_{A_0}A_1\to A_2)$.
In particular, the kernel of the canonical map $T_{A_0}(A_1)\onto A$ is generated by its degree $2$ component.
Quadraticity depends on the grading: for $n\geq2$, the grading $\deg e_{ij}=2(j-i)$ on $\Bbbk A_n$ is not quadratic, since it has no degree $1$ generators.

\begin{example}\label{ex:tensor-An}
Let $S=\bigoplus_{i=1}^n\Bbbk e_{ii}$ and $V=\bigoplus_{i=1}^{n-1}\Bbbk e_{i,i+1}$ inside $\Bbbk A_n$.
Then we have $T_S(V)\cong\Bbbk A_n$ given by\[
e_{i,i+1}\otimes_S e_{i+1,i+2}\otimes_S\cdots\otimes_S e_{j-1,j}
\longmapsto e_{ij}
\quad(1\leq i<j\leq n).
\]
\end{example}

\begin{example}\label{ex:zero-composition}
For $S,V,a,b$ from Example~\ref{ex:tensor-A3}, the quotient
\[
T_S(V)/(ab)\cong\Bbbk A_3/(e_{13})
\]
is quadratic with $R=\Bbbk(a\otimes b)$.
Its basis is $e_1,e_2,e_3,a,b$, in degrees $0,0,0,1,1$, as in Example~\ref{ex:quadratic-matrix-quotient}.
\end{example}

\begin{example}\label{ex:cubic-quotient}
Let $B=\Bbbk A_4$ and $A=B/(e_{14})$, with the grading of Example~\ref{ex:triangular-grading}.
The ideal $I=(e_{14})=\Bbbk e_{14}$ has $I_2=0$ but $I_3\neq 0$.
Since $e_{14}=e_{12}e_{23}e_{34}$, the algebra $A$ is generated in degree $1$ with a relation of degree $3$ and no relations of degree $2$.
By Proposition~\ref{prop:quadratic-basic}, any quadratic presentation would have the same generators and zero quadratic relation space, hence would give $\Bbbk A_4$ itself. Thus $A$ is not quadratic.
\end{example}

\begin{example}
The symmetric and exterior algebras of Example~\ref{ex:polynomial-exterior} are quadratic: the defining relations in \eqref{eq:polynomial-exterior} belong to $V\otimes_\Bbbk V$ and $W\otimes_\Bbbk W$, respectively.
\end{example}

\subsection{Hilbert series}

\comm{TODO: throughout, use $n$ instead of $r$ for number of simples?  Avoid using $P_i$, and just stick with $e_iA$.}

Let $A$ be a locally finite-dimensional graded algebra with $S=A_0$ finite-dimensional semisimple.
Then each degree piece of a graded right $A$-module is semisimple as $A$-module.  We call the simple direct summands in a direct sum decomposition of $M_d$  \defn{composition factors}.

\begin{notation}
By Theorem~\ref{thm:artin-wedderburn}, choose primitive idempotents $e_1,\ldots,e_r\in S$.  Then  
\[L_j:=e_jS\cong e_j A/e_jA_+\]
gives the full list of isomorphism class of (graded) simple right $S$-modules.
\phantomsection\label{not:hilbert-simples-summands}
\end{notation}

\begin{definition}\label{def:dimension-vector}\label{def:hilbert-series}
For a locally finite-dimensional graded right $A$-module $M$ with primitive idempotent $e_1, \ldots e_r$, write $[M_d:L_j]$ for the multiplicity of the composition factor $L_j$ in $M_d$.
Its degree $d$ \defn{dimension vector} is the row vector\footnote{In most representation theory literature, it is common to regard dimension vector as a column vector.  But in the context of Hilbert series, we follow Beilinson--Ginzburg--Soergel convention to use row vector here.}
\[
\dimvec M_d:=\bigl([M_d:L_j]\bigr)_{1\leq j\leq r}.
\]
The module $M$ is \defn{bounded below} if $M_d=0$ for all sufficiently negative $d$.
For such $M$, its \defn{Hilbert series} is
\[
h_M(t):=\sum_{d\in\Z}(\dimvec M_d)t^d
=\sum_{d\in\Z}\bigl([M_d:L_j]\bigr)_{1\leq j\leq r}t^d.
\]
The \defn{Hilbert series} of the algebra $A$ collects the right-module series $h_{e_iA}(t)$ as rows:
\[
h_A(t)_{ij}:=\bigl(h_{e_iA}(t)\bigr)_j
=\sum_{d\geq0}[e_iA_d:L_j]t^d
\quad\text{for all }1\leq i,j\leq r.
\]
\end{definition}

We think of the Hilbert series as ``graded dimension matrix" with one indecomposable per row. 
The entries of the Hilbert series are formal Laurent series; for algebras they are formal power series.
Finite support gives Laurent polynomials or polynomials, respectively.
Since $e_iA_0=L_i$, we have $h_A(0)=I_r$, the $r\times r$ identity array.
For a finite-dimensional right $S$-module $N$, use the same multiplicity notation and write $\dimvec N=([N:L_j])_{1\leq j\leq r}$.

The multiplicities can be computed using Hom spaces:
\begin{equation}\label{eq:composition-Hom}
[M_d:L_j]=\dim_{\End_A(L_j)}\Hom_A(e_jA,M_d)
\quad\text{for all }d\in\Z,\ 1\leq j\leq r.
\end{equation}
Indeed, $M_dA_+=0$, so every map $e_jA\to M_d$ factors uniquely through $e_jA/e_jA_+\cong L_j$.
This identifies $\Hom_A(e_jA,M_d)$ with $\Hom_A(L_j,M_d)$; precomposition gives its right vector-space structure over the division algebra $\End_A(L_j)$.
By the Schur's Lemma, A non-zero map between simple modules is an isomorphism, so each copy of $L_j$ contributes one copy of $\End_A(L_j)$, and the other simples contribute zero.
This also shows that the multiplicities are independent of the chosen decomposition.

\begin{example}\label{ex:dimension-vector-coordinate}
If $S=\bigoplus_{j=1}^r\Bbbk e_j$ with coordinate idempotents, then $\End_A(L_j)=\Bbbk$.
The Yoneda lemma~\ref{lem:idempotent-hom} says that 
\[
[M_d:L_j]=\dim_\Bbbk(M_de_j),
\qquad \dimvec M_d=(\dim_\Bbbk M_de_j)_{1\leq j\leq r}.
\]
\end{example}

\begin{lemma}\label{lem:hilbert-operations}
For locally finite-dimensional graded right $A$-modules bounded below:
\begin{enumerate}
\item A graded isomorphism $M\cong N$ implies $h_M(t)=h_N(t)$.
\item A short exact sequence $0\to L\to M\to N\to0$ gives $h_M(t)=h_L(t)+h_N(t)$.
\item For every $a\in\Z$, $h_{M(a)}(t)=t^{-a}h_M(t)$.
\end{enumerate}
\end{lemma}

\begin{proof}
A graded isomorphism preserves each degree and its simple-module multiplicities.
In each degree, a graded short exact sequence gives $0\to L_d\to M_d\to N_d\to0$, so $[M_d:L_j]=[L_d:L_j]+[N_d:L_j]$ for all $1\leq j\leq r$.
Finally, $M(a)_d=M_{d+a}$ gives
\[
h_{M(a)}(t)=\sum_d(\dimvec M_{d+a})t^d
=t^{-a}\sum_j(\dimvec M_j)t^j.\qedhere
\]
\end{proof}

\comm{From here}

\begin{definition}\label{def:total-tensor-grading}
For graded right and left $S$-modules $M,N$, their tensor product has \defn{total grading}
\[
(M\otimes_S N)_d=\bigoplus_{p+q=d}M_p\otimes_S N_q.
\]
The same definition applies to bimodules, and to vector spaces by taking $S=\Bbbk$.
\end{definition}

The decomposition follows by distributing tensor products over direct sums, as in Lemma~\ref{lem:tensor-properties}(4).
When $M,N$ are locally finite-dimensional and bounded below, only finitely many pairs $(p,q)$ contribute in a fixed degree, so the tensor product is locally finite-dimensional and bounded below.

For a finite-dimensional graded bimodule $ {}_SV_S$, use $V^{\grstar}$ and $ {}^{\grstar}V$ in Proposition~\ref{prop:S-evaluation-coevaluation}.
With the total grading, evaluation and coevaluation preserve degree: evaluation pairs degrees $i$ and $-i$, and $\id_V$ has degree $0$.

To multiply two arrays $C=(c_{ij})$ and $D=(d_{ij})$, or a row vector $m=(m_i)$ and an array $C$, use
\[
(CD)_{i\ell}=\sum_{j=1}^r c_{ij}d_{j\ell},
\qquad (mC)_j=\sum_{i=1}^r m_i c_{ij}.
\]
Thus multiplication sums over the matching middle index, rather than multiplying individual entries independently.
For series, use the same products on coefficients and collect terms with the same total power of $t$.

\begin{lemma}\label{lem:dimension-tensor}\label{lem:hilbert-tensor}
Let $N$ be a finite-dimensional right $S$-module and ${}_SU_S$ a finite-dimensional bimodule.
Then
\[
[N\otimes_S U:L_j]
=\sum_{i=1}^r[N:L_i][e_iU:L_j]
\quad\text{for all }1\leq j\leq r.
\]
For a locally finite-dimensional graded right $A$-module $M$ bounded below,
\[
h_{M\otimes_S A}(t)=h_M(t)h_A(t).
\]
\end{lemma}

\begin{proof}
Decompose $N\cong\bigoplus_{i=1}^r(e_iS)^{\oplus[N:L_i]}$.
Multiplication $e_iS\otimes_S U\to e_iU$, $e_is\otimes u\mapsto e_isu$, is an isomorphism with inverse $u\mapsto e_i\otimes u$.
Distribution over direct sums therefore gives
\[
N\otimes_S U\cong\bigoplus_{i=1}^r(e_iU)^{\oplus[N:L_i]}.
\]
Counting copies of $L_j$ proves the first formula.
Apply the first formula to $M_p\otimes_S A_q$ and sum over all $p,q\in\Z$ with $p+q=d$.
Since $e_iA_q=(P_i)_q$, the resulting coefficient is the coefficient of $t^d$ in $h_M(t)h_A(t)$.
\end{proof}
\comm{up to here.  make as supplementary}

\begin{example}[Hilbert series of a tensor algebra]\label{ex:hilbert-tensor-algebra}
Consider the tensor algebra $A=T_S(V)$ and put 
\[C:=([e_iA_1:L_j])_{1\leq i,j\leq r}=([e_iV:L_j])_{1\leq i,j\leq r}.\]
Since $e_iA_0=L_i$ and $e_iA_{d+1}\cong(e_iA_d)\otimes_S A_1$, Lemma~\ref{lem:dimension-tensor} gives, by induction,
\[
[e_iA_d:L_j]=(C^d)_{ij},
\qquad C^0=I_r
\quad\text{for all }d\geq0,\ 1\leq i,j\leq r.
\]
Consequently,
\begin{equation}\label{eq:hilbert-tensor-algebra}
\boxed{\displaystyle
h_{T_S(V)}(t)=\sum_{d\geq0}t^dC^d=(I_r-tC)^{-1}.}
\end{equation}
Indeed, multiplying the series on either side by $I_r-tC$ gives constant coefficient $I_r$; the coefficient of $t^d$ for $d\geq1$ is $C^d-C^d=0$.
The inverse is therefore an identity of formal power series with the coefficient multiplication defined above.

When $S=\Mat_n(\Bbbk)$ and $\dim_\Bbbk V=m$, the formula specialises to $h_{T_\Bbbk(V)}(t)=1/(1-mt)$.
\end{example}

\comm{make an exercise to show that if $A$ is generated in degree 1, then $h_A=h_{T_S(V)}$ yields isom $A\cong T_S(V)$, and another part of the exercise show an example when this is false if there is generator in degree 2.  put this exercise into homework 1 Q4 after part (1) before the current (2)}

\begin{example}[Upper triangular algebras and a quotient]\label{ex:hilbert-triangular}
For $A=\Bbbk A_n$, order the coordinate idempotents as $e_{11},\ldots,e_{nn}$.
Let $C$ have $C_{i,i+1}=1$ for $1\leq i<n$ and all other entries zero.
The tensor presentation in Example~\ref{ex:tensor-An} and \eqref{eq:hilbert-tensor-algebra} give
\[
h_A(t)=I_n+tC+\cdots+t^{n-1}C^{n-1},
\qquad h_A(t)_{ij}=\begin{cases}t^{j-i},&i\leq j,\\0,&i>j.\end{cases}
\]
For $n=3$, killing $e_{13}$ removes the degree $2$ contribution in entry $(1,3)$:
\[
h_{\Bbbk A_3}(t)=\begin{pmatrix}1&t&t^2\\0&1&t\\0&0&1\end{pmatrix},
\qquad
h_{\Bbbk A_3/(e_{13})}(t)=\begin{pmatrix}1&t&0\\0&1&t\\0&0&1\end{pmatrix}.
\]
The quotient removes $\Bbbk e_{13}$ from degree $2$, and hence removes $t^2$ from entry $(1,3)$.
\end{example}

\begin{example}\label{ex:hilbert-polynomial-exterior}
For the symmetric algebra $\SymAlg(V)$ and exterior algebra $\bigwedge W$ in Example~\ref{ex:polynomial-exterior} with $n=\dim_\Bbbk V, m=\dim_\Bbbk W$. 
Using the description of their basis we get
\[
h_{\SymAlg(V)}(t)=\sum_{d\ge 0}\binom{d+n-1}{n-1}t^d=\frac{1}{(1-t)^n},\footnote{differential $1/(1-t)=\sum_{d\ge 0} t^d$ on both side $n$ times to see this}\qquad
h_{\bigwedge W}(t)=(1+t)^m.
\]
\end{example}

\begin{supplementary}
\begin{lemma}\label{lem:hilbert-field-tensor}
For locally finite-dimensional graded vector spaces $U,W$ bounded below,
\[
h_{U\otimes_\Bbbk W}(t)=h_U(t)h_W(t).
\]
Let $A,B$ be locally finite-dimensional graded algebras with $A_0=\Bbbk^{\oplus r}$ and $B_0=\Bbbk^{\oplus s}$.
Write $e_1,\ldots,e_r$ and $f_1,\ldots,f_s$ for the coordinate idempotents.
Give $A\otimes_\Bbbk B$ the total grading and multiplication
$(a\otimes b)(a'\otimes b')=aa'\otimes bb'$.
Then
\[
h_{A\otimes_\Bbbk B}(t)_{(i,p),(j,q)}
=h_A(t)_{ij}h_B(t)_{pq}
\quad(1\leq i,j\leq r,\ 1\leq p,q\leq s).
\]
\end{lemma}
\end{supplementary}

\begin{proof}
For vector spaces, take dimensions in each summand of the total grading.
For algebras,
\[
(e_i\otimes f_p)(A_a\otimes_\Bbbk B_b)(e_j\otimes f_q)
=(e_iA_ae_j)\otimes_\Bbbk(f_pB_bf_q).
\]
Take vector space dimensions and sum over $a+b=d$ to obtain the coefficient of $t^d$.
\end{proof}

\begin{advanced}[suppressfirstparskip=true,innertopmargin=0pt]
\subsection{The degree 2 relations of a quotient}

Let $S$ be finite-dimensional semisimple, $V$ a finite-dimensional $S$-bimodule, and $I\subset T_S(V)$ a homogeneous two-sided ideal with $I_0=I_1=0$.
Set $A=T_S(V)/I$.
In degree $2$,
\[
A_2\cong(V\otimes_S V)/I_2,
\quad\text{where}\quad
I_2=I\cap(V\otimes_S V).
\]

\begin{definition}\label{def:quadratic-part}
The \defn{quadratic part} of $A$ is $qA:=T_S(V)/(I_2)$.
\end{definition}

Since $(I_2)\subset I$, there is a surjective graded algebra homomorphism
\[
\pi_A:qA\longrightarrow A,\qquad z+(I_2)\longmapsto z+I,
\qquad \Ker\pi_A=I/(I_2).
\]
Consequently,
\[
\pi_A\text{ is an isomorphism}
\quad\Longleftrightarrow\quad I=(I_2)
\quad\Longleftrightarrow\quad A\text{ is quadratic}.
\]
For the last equivalence, $(I_2)$ gives a quadratic presentation, while Proposition~\ref{prop:quadratic-basic} identifies the generators and degree $2$ relations of any quadratic presentation with $A_1$ and $I_2$.

\begin{example}
For $A=\Bbbk A_4/(e_{14})$ from Example~\ref{ex:cubic-quotient}, we have $I_2=0$, so
\[
qA=\Bbbk A_4,\qquad
\pi_A:\Bbbk A_4\onto\Bbbk A_4/(e_{14}),\qquad
\Ker\pi_A=\Bbbk e_{14}.
\]
Thus $qA$ does not impose the cubic relation $e_{12}e_{23}e_{34}=0$.
\end{example}
\end{advanced}

\subsection{Orthogonal complements and quadratic duality}

Let $S$ be finite-dimensional semisimple.

\begin{definition}\label{def:orthogonals}
For an $S$-sub-bimodule $U$ of a finite-dimensional $S$-bimodule $M$, its \defn{right orthogonal complement} and \defn{left orthogonal complement} are
\[
\begin{aligned}
U^\perp&:=\{f\in M^*\mid f(u)=0\text{ for all }u\in U\},\\
{}^\perp U&:=\{g\in{}^*M\mid g(u)=0\text{ for all }u\in U\}.
\end{aligned}
\]
\end{definition}

Both orthogonal complements are $S$-sub-bimodules.
Indeed, for $s,t\in S$ and $u\in U$, \eqref{eq:tensor-dual-actions} gives
\[
(sft)(u)=sf(tu)=0\quad(f\in U^\perp),\qquad
(sgt)(u)=g(us)t=0\quad(g\in{}^\perp U),
\]
since $tu,us\in U$.
The orthogonal complements are defined without choosing a symmetrising form, but it is possible to use realise (through the natural isomorphism associated to a chosen symmeterising form) the orthogonal complements as a sub-bimodule of $\bbD(M)$.

\begin{advanced}
If $S$ is equipped with a symmetrising form $\tau:S\to\Bbbk$, the comparison maps $\alpha_M,\beta_M$ of \eqref{eq:dual-comparison} satisfy
\begin{equation}\label{eq:orthogonal-comparison}
\alpha_M(U^\perp)=\beta_M({}^\perp U)
=\{\lambda\in D(M)\mid\lambda(u)=0\text{ for all }u\in U\}.
\end{equation}
The forward inclusions follow by composing with $\tau$.
Conversely, $\alpha_M(f)(U)=0$ gives
\[
\tau(f(u)s)=\tau(f(us))=0
\quad\text{for all }s\in S,\ u\in U.
\]
Non-degeneracy of $\tau$ forces $f(u)=0$ for every $u\in U$.
For the left dual, $\tau(sg(u))=\tau(g(su))=0$ gives $g(u)=0$ instead.
Bijectivity of $\alpha_M,\beta_M$ in Theorem~\ref{prop:dual-comparison} proves \eqref{eq:orthogonal-comparison}.
\end{advanced}

\begin{example}\label{ex:symmetric-orthogonals}
We compute the orthogonal complements of the relation spaces of $\SymAlg(V)$ and $\bigwedge W$ for $\dim_{\Bbbk}V=2=\dim_{\Bbbk}W$.
Take $V=\Bbbk x_1\oplus\Bbbk x_2$ and $W=\Bbbk y_1\oplus\Bbbk y_2$ in Example~\ref{ex:polynomial-exterior}.
The relation spaces in $V\otimes_\Bbbk V$ and $W\otimes_\Bbbk W$, respectively, are
\[
\begin{aligned}
R_{\SymAlg(V)}&=\Bbbk(x_1\otimes x_2-x_2\otimes x_1),\\
R_{\bigwedge W}&=\Bbbk(y_1\otimes y_1)\oplus\Bbbk(y_2\otimes y_2)\oplus\Bbbk(y_1\otimes y_2+y_2\otimes y_1).
\end{aligned}
\]
Let $x_1^*,x_2^*$ and $y_1^*,y_2^*$ be the dual bases of $D(V)$ and $D(W)$, respectively.
Under the tensor-dual identification \eqref{eq:tensor-evaluations},
\[
\begin{aligned}
R_{\SymAlg(V)}^\perp&=\operatorname{span}_\Bbbk\{x_1^*\otimes x_1^*,\,x_2^*\otimes x_2^*,\,x_1^*\otimes x_2^*+x_2^*\otimes x_1^*\},\\
R_{\bigwedge W}^\perp&=\Bbbk(y_1^*\otimes y_2^*-y_2^*\otimes y_1^*).
\end{aligned}
\]
The first space consists of the exterior relations on $D(V)$; the second consists of the symmetric relations on $D(W)$.
\end{example}

\Needspace{12\baselineskip}
\subsection{Exercises}

\begin{exercise}[Interaction of tensor product with $S$-dual]\label{ex:tensor-duality-proof}
Let $S$ be finite-dimensional semisimple, and consider finite-dimensional bimodules ${}_SV_S$, ${}_SW_S$.
Prove the right $S$-dual isomorphism \eqref{eq:tensor-duals}, the semisimple case of Proposition~\ref{prop:tensor-duals}.
\begin{enumerate}
\item Check that the formula for $\theta_r$ in \eqref{eq:tensor-evaluations} is bilinear in each pair of variables, is unchanged by
\[
vs\otimes w\longleftrightarrow v\otimes sw,
\qquad gs\otimes f\longleftrightarrow g\otimes sf
\quad(s\in S),
\]
and gives a right $S$-linear map $V\otimes_S W\to S$ for each $g\otimes f$.
Verify $\theta_r(sg\otimes ft)=s\theta_r(g\otimes f)t$ for $s,t\in S$.
\item Prove that $\theta_r$ is bijective.
\item For bimodule maps $p:V\to V'$ and $q:W\to W'$, prove
\[
(p\otimes q)^*\circ\theta_r^{V',W'}
=\theta_r^{V,W}\circ(q^*\otimes p^*).
\]
Here the superscripts specify the pair of bimodules, and $p^*,q^*$, and $(p\otimes q)^*$ are given by precomposition with $p,q$, and $p\otimes q$, respectively.
\end{enumerate}
\emph{Hint for (2).} First take $V=S$, then use Corollary~\ref{cor:semisimple-summands}.

\emph{Remark.} The analogous assertions hold for left $S$-duals, with left and right interchanged.
\end{exercise}

\begin{hidden}
\begin{proof}
\emph{Well-definedness.}
For the right duals, right $S$-linearity gives
\[
g(f(vs)w)=g(f(v)sw),\qquad
(gs)(f(v)w)=g(sf(v)w)=g((sf)(v)w).
\]
The first identity respects the relation in $V\otimes_S W$, and the second respects the relation in $W^*\otimes_S V^*$.
The resulting functional is right $S$-linear because $g$ is.
For the left duals, the corresponding identities are
\[
f(vs\,g(w))=f(vg(sw)),\qquad
f(v(gs)(w))=f(vg(w)s)=(sf)(vg(w)).
\]
Left $S$-linearity of $f$ makes the resulting functional left $S$-linear.
The outer actions in \eqref{eq:tensor-dual-actions} give
\[
\begin{aligned}
\theta_r(sg\otimes ft)(v\otimes w)&=s\,g(f(tv)w),\\
\theta_\ell(sg\otimes ft)(v\otimes w)&=f(vg(ws))t,
\end{aligned}
\]
which are the required bimodule identities.

\emph{Bijectivity.}
To check $\theta_r$, forget the left action on $V$.
For $V=S$, identify $S^*\cong S$ by $f\mapsto f(1)$; then $\theta_r$ is $W^*\otimes_S S\to W^*$, $g\otimes s\mapsto gs$, an isomorphism by Lemma~\ref{lem:tensor-properties}(3).
The formula commutes with maps of right $S$-modules in $V$, so it is an isomorphism for finite direct sums and direct summands of $S$.
Corollary~\ref{cor:semisimple-summands} applies to $V$.
For $\theta_\ell$, use $W=S$ as a left module, where the map is $S\otimes_S{}^*V\to{}^*V$, $s\otimes f\mapsto sf$, and apply the same argument to $W$.

\emph{Naturality.}
For bimodule maps $p:V\to V'$ and $q:W\to W'$, the square
\[
\xymatrix@C=48pt@R=22pt{
(W')^*\otimes_S(V')^* \ar[r]^{\theta_r} \ar[d]_{q^*\otimes p^*}
& (V'\otimes_S W')^* \ar[d]^{(p\otimes q)^*}\\
W^*\otimes_S V^* \ar[r]_{\theta_r}
& (V\otimes_S W)^*
}
\]
commutes: both routes send $g\otimes f$, evaluated at $v\otimes w$, to $g(f(p(v))q(w))$.
For $\theta_\ell$, the corresponding two routes give $f(p(v)g(q(w)))$.
\end{proof}

\end{hidden}

\begin{exercise}
Let $A=\Bbbk A_4/(e_{13},e_{24})$ with the grading of Example~\ref{ex:triangular-grading}.
Find a homogeneous basis, compute $h_A(t)$, and write $A=T_S(V)/(R)$ with explicit $S,V,R$.
\emph{Hint.} Use Lemma~\ref{lem:idempotent-tensors}; the ideal contains $e_{14}$ as well as $e_{13},e_{24}$.
\end{exercise}

\begin{exercise}
For the bimodule $V$ in Example~\ref{ex:two-vertex-relations}, find a basis of each $T_S(V)_d$ and show that
\[
h_{T_S(V)}(t)=\frac{1}{1-t^2}\begin{pmatrix}1&t\\t&1\end{pmatrix}.
\]
Compare with the quotient $T_S(V)/(ab)$.
\emph{Hint.} Count alternating words separately in each degree, including degree $0$.
\end{exercise}

\begin{exercise}
For the algebras in Example~\ref{ex:polynomial-exterior}, prove
\[
\dim_\Bbbk\SymAlg(V)_2=\frac{n(n+1)}2,\qquad
\dim_\Bbbk(\bigwedge W)_2=\frac{m(m-1)}2
\quad\text{where }n=\dim_\Bbbk V,\ m=\dim_\Bbbk W.
\]
\emph{Hint.} Count the ordered pairs allowed in the two bases of Example~\ref{ex:polynomial-exterior}.
\end{exercise}

\pagebreak

